\chapter*{ВВЕДЕНИЕ}
\addcontentsline{toc}{chapter}{ВВЕДЕНИЕ}
\sloppy

За последние несколько лет глубокое обучение пришло в геномику в форме языковых моделей ДНК (DNA Language Models, DNA~LM). По сути, идея прямая: берём нуклеотидные последовательности, обучаем на них модель в режиме самосупервизии и получаем представления, полезные для задач вроде предсказания регуляторных элементов, эффектов мутаций и аннотации генома. HyenaDNA~\cite{nguenHyenaDNALongRangeGenomic2024}, DNABERT-2~\cite{zhouDNABERT2EfficientFoundation2024}, Caduceus~\cite{schiffCaduceusBiDirectionalEquivariant2024}, Nucleotide Transformer~\cite{dalla-torreNucleotideTransformerBuilding2025}, Evo~\cite{nguyenSequenceModelingDesign2024} - все эти модели показывают, что подход действительно работает~\cite{benegasGenomicLanguageModels2024}.

Однако в этой картине есть узкое место, которому до недавнего времени уделялось мало внимания: токенизация. Прежде чем модель увидит последовательность, её нужно разбить на дискретные единицы - токены. В геномике обычно используют три стандартных подхода: посимвольную обработку (каждый нуклеотид A, T, C, G - отдельный токен), фиксированные $k$-меры и BPE~\cite{sennrichNeuralMachineTranslation2016}, заимствованный из NLP~\cite{tokenizationGenomicsReview2025}. При этом все три подхода статические: правила разбиения фиксируются до обучения и не зависят от того, что именно закодировано в последовательности.

Проблема здесь в том, что геном неоднороден. Регуляторные элементы - промоторы, энхансеры, сайты связывания транскрипционных факторов - сосредоточены на коротких участках и требуют высокого разрешения. В то же время межгенные области и повторы занимают большую часть генома и несут меньше функциональной информации~\cite{schoenfelderLongrangeEnhancerPromoter2019}. Фиксированные $k$-меры разрывают функциональные мотивы на произвольных границах. BPE, в свою очередь, оптимизирует сжатие по частотности подстрок, а не по биологической значимости. Поэтому повторяющиеся элементы (Alu, LINE), составляющие около 45\% генома~\cite{landerInitialSequencingAnalysis2001}, получают компактные токены, а уникальные регуляторные элементы - нет~\cite{lindseyImpactTokenizerSelection2025}. Исследования Vishniakov~и~др.~\cite{vishniakovGenomicFoundationlessModels2024} и Eapen~\cite{eapenGenomicTokenizerBiologydriven2025} подтверждают, что выбор токенизатора существенно влияет на качество модели, но систематическое сравнение адаптивных подходов с учётом биологических приоров пока не проведено.

На этом фоне в 2024--2026~годах появился целый ряд работ по обучаемой токенизации для ДНК: MxDNA~\cite{qiaoMxDNAAdaptiveDNA2024} на основе смеси экспертов, MergeDNA~\cite{mergeDNAContextAware2026} с контекстно-зависимым слиянием, DNAChunker~\cite{dnachunkerLearnableTokenization2026}, EvoLen~\cite{evoLenEvolutionGuided2026}, PatchDNA~\cite{patchDNA2026}, BioToken~\cite{bioTokenBioFM2025}. Параллельно в NLP развиваются подходы к адаптивной byte-level токенизации: Byte Latent Transformer~\cite{bltByteLatentTransformer2024}, MegaByte~\cite{yuMegaBytePredictingMillionbyte2023}, Charformer~\cite{tayCharformerFastCharacter2022}. Вместе эти работы формируют контекст, в котором становится возможным предложить новый подход к токенизации геномных данных.

Цель настоящей работы~- разработать метод адаптивной токенизации ДНК-последовательностей, основанный на рекурсивном гейтированном слиянии с интеграцией биологических приоров, и оценить его влияние на качество предобучения и решения задач геномики.

Для этого поставлены следующие задачи:
\begin{enumerate}
  \item проанализировать существующие методы токенизации ДНК и выявить их ограничения;
  \item разработать архитектуру адаптивного токенизатора на основе рекурсивного гейтированного слияния с дифференцируемым динамическим программированием и интеграцией биологических приоров (оценки консервативности, мотивы транскрипционных факторов);
  \item интегрировать разработанный метод в языковую модель ДНК с каузальным декодером;
  \item провести экспериментальную оценку на бенчмарке DART-Eval~\cite{patel2025dartevalcomprehensivednalanguage} в сравнении с существующими методами.
\end{enumerate}

Объект исследования~- языковые модели для нуклеотидных последовательностей ДНК.

Предмет исследования~- методы адаптивной токенизации ДНК-последовательностей для повышения качества геномных языковых моделей.

Научная новизна работы состоит в следующем:
\begin{enumerate}
  \item предложена архитектура адаптивного токенизатора на основе рекурсивного гейтированного слияния с дифференцируемым DP через log-semiring relaxation~\cite{menschDifferentiableDynamicProgramming2018}, которая позволяет оптимизировать границы токенов непрерывно, сохраняя ограничение непересечения;
  \item разработан механизм интеграции биологических приоров (оценки эволюционной консервативности phastCons~\cite{siepelEvolutionarilyConservedElements2005} и позиционные весовые матрицы мотивов JASPAR~\cite{rauluseviciuteJASPAR2024}) в функцию оценки слияния;
  \item предложена составная функция потерь, объединяющая NTP с регуляризаторами сжатия, энтропии решений, RC-согласованности и разнообразия длин спанов.
\end{enumerate}

Практическая значимость работы состоит в том, что разработанный токенизатор может быть встроен как модуль предобработки в существующие архитектуры геномных языковых моделей. Кроме того, реализация выполнена на Python с использованием PyTorch~\cite{paszke2019pytorchimperativestylehighperformance} и совместима с типовыми пайплайнами обучения геномных моделей.

Работа состоит из введения, двух глав, заключения, списка использованных источников и приложения. В первой главе рассматриваются существующие геномные языковые модели, методы токенизации ДНК и бенчмарки. Вторая глава описывает предлагаемый метод: архитектуру рекурсивного слияния, механизм дифференцируемого выбора пар, функцию потерь, стратегию обучения и план экспериментов.
